% ch07_u09 —— 原书 p.316--p.325（7.11.15(iii) -- 7.12.17 中段）
% 承接 ch07_u08：本单元自 7.11.15 的 (iii) 续起，不写 \bdchap。
% 尾部接 p.326（7.12.17 中（356）式之后的句子未完）。

\SourcePage{316}

(iii) 若 $F_1,F_2$ 分别是 $X_1,X_2$ 上的松弛复形，则 $F_1\boxtimes F_2$ 是 $X_1\times X_2$ 上的松弛 Ω-复形。

\emph{证明.} 按显然的方式修改 7.3.8 的证明。\hfill$\square$

由此可见，导出范畴 $D(X)$ 可以用松弛复形来定义。

\bdp{7.11.16} 归纳概形的任意态射 $f:X\to Y$ 给出推前函子
\[
f_{\cdot}:C(X,\Omega)\to C(Y,\Omega),
\]
它保持子范畴 $C_c$。$f_{\cdot}$ 把具拟紧支撑的松弛复形之间的 D-拟同构仍送到 D-拟同构，这一点留给读者核验（参看 7.3.9、7.3.11(ii)）。于是右导出函子 $Rf_{\cdot}=f_{*}:D(X)\to D(Y)$ 良定义：若 $F$ 是具拟紧支撑的松弛复形，则 $f_{*}F=f_{\cdot}F$。由于 $f_{\cdot}$ 把松弛复形送到松弛复形，$f_{*}$ 与 $f$ 的复合相容。

对 $M\in D(X,\mathcal{D})$，记 $M_{\mathcal{O}}\in D(X,\mathcal{O})$ 为把同一个 $M$ 看作 $\mathcal{O}^!$-模复形。有典范的积分态射
\[
i_f:Rf_{\cdot}(M_{\mathcal{O}})\to(f_{*}M)_{\mathcal{O}},
\]
它位于 $D(Y,\mathcal{O})$ 中，定义同 7.2.11；该态射与 $f$ 的复合相容。

\bdp{7.11.17} 按 7.6.1 定义 Hecke 幺半范畴 $\mathcal{H}$。出发点是 ind-affine 群归纳概形 $G$ 及其仿射群子概形 $K\subset G$。假设 $G/K$（关于 fpqc 拓扑的层商）是 ind-finite type 的归纳概形；它自动形式光滑，且对角态射仿射。显然 $G$ 是合理归纳概形，$K$ 是它的合理子概形。考虑 $G$ 上具拟紧支撑的 $(K\times K)_{\Omega}$-等变 $\Omega^!$-复形构成的 DG 范畴 $\mathcal{H}^c$（见 7.11.12 注 (i)）。由下降，这样的复形等同于 $G/K$ 上或 $K\backslash G$ 上的 $K_{\Omega}$-等变容许 $\Omega^!$-复形。相应的 D-拟同构概念彼此等价。我们的 $\mathcal{H}$ 就是相应的 D-导出范畴。
% [Tutor Note: 原书 p.316 此处印作 "is a ind-scheme of ind-finite type"，冠词用 a 而不用 an；底本如此，中文无对应形式。]

\SourcePage{317}

7.6.1 的构造在我们的情形下完全适用。故 $\mathcal{H}^c$ 是 DG 幺半范畴，$\mathcal{H}$ 是三角幺半范畴。

\bdp{7.11.18} 设概形 $Y$ 带 $G$-作用，且存在开拟紧子概形的递增族 $U_0\subset U_1\subset\ldots$，使 $Y=\bigcup U_i$，并使得对某个合理群子概形 $K_i\subset G$，$K_i$ 在 $U_i$ 上的作用自由，而 $K_i\backslash U_i$ 是光滑概形（特别地有限型）。于是叠 $\mathcal{B}=K\backslash Y$ 光滑（它被诸概形 $(K_i\cap K)\backslash U_i$ 覆盖），其对角态射仿射，故可用 7.3.12 的 $D(\mathcal{B})$ 定义。

为在 $D(\mathcal{B})$ 上定义 $\mathcal{H}$-作用，照 7.6.1 的办法做，并作如下修改（这些修改源于 $Y$ 与 $G$ 可能非拟紧）。可设上述 $U_i$ 都 $K$-不变；置 $\mathcal{B}_i=K\backslash U_i\subset\mathcal{B}$。取松弛 Ω-复形 $F=\bigcup F_n\in C_a(K\backslash G/K,\Omega)$（故 $F_n$ 的支撑 $S_n$ 拟紧）与 $T\in C(\mathcal{B}_{\cdot},\Omega)$。取 $j(n,i)$ 为关于 $n$、$i$ 都递增的序列，使 $S_n^{-1}\cdot U_i\subset U_{j(n,i)}$。考虑 $\mathcal{B}_i$ 上的 Ω-复形
\[
(F_n\circledast T)_i:=\bar m_{U_i}\cdot p^{\cdot}_{U_i\Omega}(F_n\boxtimes T_{j(n,i)})|_{\mathcal{B}_i}
\]
与
\[
(F_n\circledast T)'_i:=\bar m_{U_i}\cdot p^{\cdot}_{U_i\Omega}(F_n\boxtimes T_{j(n+1,i)})|_{\mathcal{B}_i}.
\]
有显然的态射 $(F_n\circledast T)'_i\to(F_{n+1}\circledast T)_i$ 与 $(F_n\circledast T)'_i\to(F_n\circledast T)_i$，后者是拟同构。置
\[
(F\circledast T)_i:=\operatorname{Cone}\big(\oplus(F_n\circledast T)_i\to\oplus(F_n\circledast T)_i\big),
\]
其中箭头是上述态射的（逐分量）差。
% [Tutor Note: 原书 p.317 此处 Cone 的两个直和项都印成 $(F_n\circledast T)_i$；按上下文第二项应为 $(F_n\circledast T)'_i$，底本如此，保留原样。]

这些 $(F\circledast T)_i$ 按显然的方式构成 $C(\mathcal{B},\Omega)$ 的对象 $F\circledast T$。留给读者核验：作为 $D(\mathcal{B})$ 的对象，$F\circledast T$ 不依赖辅助数据（$U_i$ 与 $j(n,i)$）的选取；且 $\circledast$ 是 $D(\mathcal{B})$ 上的 $\mathcal{H}$-作用。
% [Tutor Note: 原书 p.317 本节两处把「作用」印作「Action」（H-Action），无故大写；底本如此，译文照常作「作用」。]

\bdsec{7.12. 归纳概形与 Mittag-Leffler 模}
% [Tutor Note: 原书 p.317 的 7.12 节标题与正文同行（run-in 粗体）；底本如此，译文将标题单独排出。]

Raynaud 与 Gruson [RG] 引入了 Mittag-Leffler 模这一非凡的概念。本节说明：平坦 Mittag-Leffler 模这一概念，在某种意义上正是 ind-finite type 形式光滑归纳概形概念的线性化版本（见 7.12.12、7.12.14、7.12.15）。
\SourcePage{318}%
利用可数生成的平坦 Mittag-Leffler 模必投射这一事实，我们描述 ind-finite type 的形式光滑仿射 $\aleph_0$-形式概形（见 7.12.22、7.12.23）。

读者可以跳过本节：它的结果在本书其余部分用不到（我们收录它们，只为澄清形式光滑归纳概形这一概念）。

7.11 中我们假定「归纳概形」意即「$\mathbb{C}$ 上的归纳概形」（这其实无关紧要）。本节我们宁可放弃这个假设。

\bdp{7.12.1} 设 $A$ 是环\footnote{我们假设 $A$ 交换，但在 7.12.1--7.12.8 中这并不必要（只需按显然的方式在「模」字前加上「左」或「右」）。}。记 $\mathcal{C}$ 为有限表现 $A$-模的范畴。据 [RG] p.69，称 $A$-模 $M$ 为 \emph{Mittag-Leffler 模}，若任意态射 $f:F\to M$（$F\in\mathcal{C}$）都可分解为 $F\xrightarrow{u}G\to M$（$G\in\mathcal{C}$），使得对 $f$ 的任意分解 $F\xrightarrow{u'}G'\to M$（$G'\in\mathcal{C}$），存在态射 $\varphi:G'\to G$ 使 $u=\varphi u'$。

\bdp{7.12.2} 设 $M=\varinjlim M_i$，$i\in I$，其中 $I$ 是有向有序集，$M_i\in\mathcal{C}$。据 loc.cit.，$M$ 是 Mittag-Leffler 模当且仅当：对每个 $i\in I$ 存在 $j\ge i$，使得对每个 $k\ge i$，态射 $u_{ij}:M_i\to M_j$ 都可分解为 $\varphi_{ijk}u_{ik}$，其中 $\varphi_{ijk}:F_k\to F_j$。
% [Tutor Note: 原书 p.318 此处印作 $\varphi_{ijk}:F_k\to F_j$，而上文用的是 $M_i$ 记号；底本如此，保留原样。]

$I$ 是滤过范畴时也一样；若 $I$ 是「$\mathcal{C}$ 的对象到 $M$ 的全部态射」构成的范畴，而 $F_i\in\mathcal{C}$ 是态射 $i$ 的源，则上述陈述就是同义反复。

\bdp{7.12.3} 上述归纳系统 $(M_i)$ 的这一性质，在把 $\mathcal{C}$ 换成任意范畴 $\mathcal{C}'$ 后仍有意义。若 $\mathcal{C}'$ 是集合范畴的对偶，即有集合的射影系统 $(E_i,u_{ij}:E_j\to E_i)$，则由 EGA $0_{\mathrm{III}}$ 13.1.2 得到下述 Mittag-Leffler 条件：对每个 $i\in I$ 存在 $j\ge i$，使得 $u_{ij}(E_j)=u_{ik}(E_k)$ 对所有 $k\ge j$ 成立。该条件成立当且仅当射影系统 $(E_i,u_{ij})$ 等价于射影系统 $(\tilde E_{\alpha},\tilde u_{\alpha\beta})$，其中映射 $\tilde u_{\alpha\beta}$ 都是%
\SourcePage{319}%
满的。证「仅当」只需对充分大的 $j$ 置 $\tilde E_i:=u_{ij}(E_j)$。

\bdp{7.12.4} 设 $M=\varinjlim M_i$，$M_i\in\mathcal{C}$。据 [RG]，$M$ 是 Mittag-Leffler 模当且仅当：对任意从 $\mathcal{C}$ 到集合范畴的反变函子 $\Phi$，射影系统 $(\Phi(M_i))$ 都满足 Mittag-Leffler 条件（证「若」时考虑函子 $\Phi(N)=\operatorname{Hom}(N,\prod_i M_i)$ 或 $\tilde\Phi(N)=\coprod_i\operatorname{Hom}(N,M_i)$）。设 $M$ 平坦，置 $M^{*}_i=\operatorname{Hom}(M_i,A)$。据 [RG]，$M$ 是 Mittag-Leffler 模当且仅当射影系统 $(M^{*}_i)$ 满足 Mittag-Leffler 条件。$M_i$ 投射时这显然。一般情形由 Lazard 引理得出（存在由有限生成投射模构成、与 $(M_i)$ 等价的归纳系统）。

\bdp{7.12.5} 考虑从 $A$-模范畴到阿贝尔群范畴的两类函子：

1) 对 $A$-模 $M$ 有函子
\begin{equation}\tag{354}
L\mapsto L\otimes_A M;
\end{equation}

2) 对 $A$-模的射影系统 $N_i$（$i$ 属于某个有向有序集）有函子
\begin{equation}\tag{355}
L\mapsto\varinjlim_i\operatorname{Hom}(N_i,L).
\end{equation}
% [Tutor Note: 原书 p.319 的（354）（355）与 p.325 的（356）编号都排在左缘；底本如此，译文按通行的右缘排法排印。]

\bdp{7.12.6} \emph{命题.}
\begin{enumerate}[label=(\roman*),leftmargin=2.2em]
\item 函子 (354) 同构于形如 (355) 的函子，当且仅当 $M$ 平坦；
\item 函子 (354) 同构于射影系统 $(N_i)$ 对应的函子 (355)（转移映射 $N_j\to N_i$，$i\le j$，都满），当且仅当 $M$ 是平坦 Mittag-Leffler 模；
\item 对应满转移映射射影系统的函子 (355) 同构于形如 (354) 的函子，%
\SourcePage{320}%
当且仅当 (355) 正合且各 $N_i$ 有限生成。
\end{enumerate}

\emph{证明.} 若 (354) 与 (355) 同构，则 (354) 左正合，故 $M$ 平坦。若 $M$ 平坦，则由 Lazard 引理 $M=\varinjlim P_i$，其中 $P_i$ 投射且有限生成；于是对应 $N_i=P^{*}_i$ 的 (355) 同构于 (354)。(i) 得证。由 (i) 推出 (ii)：对上述 $P_i$，射影系统 $(P^{*}_i)$ 等价于转移映射满的 $(N_i)$，当且仅当 $(P^{*}_i)$ 满足 Mittag-Leffler 条件（见 7.12.3）。证 (iii)：形如 (354) 的函子正是右正合且与无限直和交换（从而与归纳极限交换）的加性函子。形如 (355) 的函子右正合当且仅当它正合。若各 $N_i$ 有限生成，则 (355) 与无限直和交换。若转移映射满且 (355) 与归纳极限交换，则各 $N_i$ 有限生成。\hfill$\square$

\bdp{7.12.7} 由 7.12.6，平坦 Mittag-Leffler 模「等同于」有限生成模的射影系统 $(N_i)$ 的等价类（转移映射 $N_j\to N_i$（$i\le j$）都满，且函子 (355) 正合）。更确切地说，$M=\varinjlim_i\operatorname{Hom}(N_i,A)$（此时 (354) 与 (355) 同构）。

\bdp{7.12.8} \emph{定理}(Raynaud--Gruson).（D. Lazard 呢？据 [RG] p.73，此想法可溯至 D. Lazard 学位论文第一章的定理 3.1 与 3.2，Bull. Soc. Math. France 第 97 卷 (1969), 81--128；另见他在同刊第 95 卷 (1967), 95--108 的工作。）下列条件等价：
\begin{enumerate}[label=(\roman*),leftmargin=2.2em]
\item $M$ 是平坦 Mittag-Leffler 模；
\SourcePage{321}%
\item $M$ 的每个有限或可数子集都含于某个可数生成的投射子模 $P\subset M$，且 $M/P$ 平坦；
\item $M$ 的每个有限子集都含于某个投射子模 $P\subset M$，且 $M/P$ 平坦。
\end{enumerate}

特别地，投射模都是 Mittag-Leffler 模，而可数生成\footnote{可数生成的假设不可少；见 7.12.24。}的平坦 Mittag-Leffler 模都是投射模。

(iii)$\Rightarrow$(i) 容易。（只须证：若 $F,F'$ 有限表现，$\varphi:F\to F'$，$\psi:F'\to M$，且 $\psi\varphi(F)\subset P$，则存在 $\tilde\psi:F'\to M$ 使 $\tilde\psi(F')\subset P$ 且 $\tilde\psi\varphi=\psi\varphi$；这要用到 $\operatorname{Hom}(L,M)\to\operatorname{Hom}(L,M/P)$ 对有限表现 $L$ 满，特别取 $L=\operatorname{Coker}\varphi$。）

(i)$\Rightarrow$(ii) 见 [RG] p.73--74。关键论证如下。设序列 $P_1\to P_2\to\ldots$，各 $P_i$ 有限生成投射，且射影系统 $(P^{*}_i)$ 满足 Mittag-Leffler 性质。要证 $P:=\varinjlim P_i$ 投射，只须证：对任意正合列 $0\to N'\to N\to N''\to 0$，映射 $\operatorname{Hom}(P,N)\to\operatorname{Hom}(P,N'')$ 满。对每个 $i$，序列 $0\to P^{*}_i\otimes N'\to P^{*}_i\otimes N\to P^{*}_i\otimes N''\to 0$ 正合，问题归结为这些序列的射影极限正合。据 EGA $0_{\mathrm{III}}$ 13.2.2，这由射影系统 $(P^{*}_i\otimes N')$ 的 Mittag-Leffler 性质推出。

\emph{注.} 若指标集 $i$ 不可数，我们就不能\footnote{EGA $0_{\mathrm{III}}$ 13.2.2 的论证基于下述事实：由可数集 $I$ 参数化的非空集射影系统 $(Y_i)_{i\in I}$，若满足 Mittag-Leffler 条件，则其射影极限非空。不可数情形下这不对。例如取不可数集 $S$，对每个有限 $F\subset A$ 记 $Y_F$ 为单射 $F\to\mathbb{N}$ 的集合；映射 $Y_{F'}\to Y_F$（$F'\supset F$）都满，但 $\varprojlim_F Y_F=\emptyset$。}用 EGA $0_{\mathrm{III}}$ 13.2.2。
% [Tutor Note: 原书 p.321 脚注前文设的是不可数集 $S$，此处印作「有限 $F\subset A$」，$A$ 应为 $S$；底本如此，保留原样。]

\SourcePage{322}

$P$ 投射性的另一证明如下（其实是同一个证明的另一版本）。记 $f_i:P_i\to P_{i+1}$。Mittag-Leffler 性质说，把序列 $\{P_i\}$ 换成它的子序列后，存在 $g_i:P_{i+1}\to P_i$ 使 $g_{i+1}f_{i+1}f_i=f_i$。置 $\mathcal{P}:=\oplus_i P_i$。记 $f:\mathcal{P}\to\mathcal{P}$、$g:\mathcal{P}\to\mathcal{P}$ 为 $f_i$、$g_i$ 诱导的算子。则 $gf^2=f$。有正合列
\[
0\to\mathcal{P}\xrightarrow{1-f}\mathcal{P}\to P\to 0.
\]
因 $\mathcal{P}$ 投射，只须证该列分裂，即有 $h:\mathcal{P}\to\mathcal{P}$ 使 $h(1-f)=1$。事实上置 $h=1-(1-g)^{-1}gf$，并用等式 $gf^2=f$。\footnote{D. Arinkin 注意到，先验地显然：若 $f,g$ 是一个（非交换）环 $R$ 的元素，$gf^2=f$，且 $1-g$ 有左逆，则 $1-f$ 有左逆。事实上，记 $\mathbf{1}$ 为 $1$ 在 $R/R(1-f)$ 中的像，则 $f\mathbf{1}=\mathbf{1}$、$gf^2\mathbf{1}=g\mathbf{1}$，故 $g\mathbf{1}=\mathbf{1}$，从而 $\mathbf{1}=0$。}

\bdp{7.12.9} \emph{命题.} 设 $B$ 是 $A$-代数。若 $M$ 是 Mittag-Leffler $A$-模，则 $B\otimes_A M$ 是 Mittag-Leffler $B$-模。若 $B$ 在 $A$ 上忠实平坦，则反之亦然。

这在 [RG] 中有证明。证明容易：把 $M$ 表为有限表现模的归纳极限，再用 7.12.2。

于是，概形上 Mittag-Leffler $\mathcal{O}$-模的概念、归纳概形上 Mittag-Leffler $\mathcal{O}^p$-模的概念，都清楚了。

\bdp{7.12.10} \emph{命题.} 诺特概形 $S$ 上可数型的平坦 Mittag-Leffler $\mathcal{O}$-模 $\mathcal{F}$ 局部自由。若 $S$ 仿射且连通，而 $\mathcal{F}$ 无限型，则 $\mathcal{F}$ 自由。

这是 7.12.8 与下述结果的直接推论。

\bdp{7.12.11} \emph{定理.} 若 $R$ 诺特且 $\operatorname{Spec}R$ 连通，则每个不是有限生成的投射 $R$-模都自由。

本定理由 Bass 证明（见 [Ba63] 的推论 4.5）。

\SourcePage{323}

\bdp{7.12.12} \emph{命题.} 设 $X$ 是域上 ind-finite type 的形式光滑归纳概形。则 $\mathcal{O}^p$-模 $\Theta_X$、$\mathcal{D}_X$、$\mathcal{D}_{iX}$（见 7.11.8）都是平坦 Mittag-Leffler 模。

\emph{证明.} 我们先证 $\mathcal{D}_X$ 到闭子概形 $Y\subset X$ 的限制是平坦 Mittag-Leffler $\mathcal{O}_Y$-模（对 $\Theta_X$ 与 $\mathcal{D}_{iX}$ 的论证相同）。可设 $Y$ 仿射（否则用 $X\setminus F$ 代替 $X$，其中 $F\subset Y$ 适当闭）。据 7.12.6，只须证：
\begin{enumerate}[label=(\roman*),leftmargin=2.2em]
\item $\mathcal{O}_Y$-模范畴上定义的函子 $L\mapsto L\otimes\mathcal{D}_X$ 正合；
\item 它具有 (355) 的形式，其中 $\mathcal{O}_Y$-模 $N_i$ 凝聚。
\end{enumerate}
按定义，$L\otimes\mathcal{D}_X$ 就是 (349) 所定义的层 $\mathcal{D}(L)$，故 (ii) 显然。(i) 我们已在 7.11.8 证明。\hfill$\square$

\bdp{7.12.13} \emph{命题.} 设 $X$ 是域上 ind-finite type 的形式光滑 $\aleph_0$-归纳概形，$Y\subset X$ 是局部闭子概形。则 $\Theta_X$ 到 $Y$ 的限制局部自由。若 $Y$ 仿射且连通，而 $\Theta_X$ 到 $Y$ 的限制无限型，则它自由。

这由 7.12.12 与 7.12.10 推出。

\bdp{7.12.14} \emph{命题.} 设 $A$ 是环，$M$ 是 $A$-模。由 $F_M(R)=M\otimes R$ 定义「$A$-空间」$F_M$（即从 $A$-代数范畴到集合范畴的函子）。则 $F_M$ 是归纳概形当且仅当 $M$ 是平坦 Mittag-Leffler 模；此时 $F_M$ 在 $A$ 上形式光滑，且在 $A$ 上 ind-finite type。

\emph{证明.} 若 $M$ 是平坦 Mittag-Leffler 模，则由 7.12.6(ii) $F_M$ 是归纳概形，由 7.12.6(iii) 它在 $A$ 上 ind-finite type。形式光滑由定义得出。反之，设 $F_M$ 是归纳概形。把 $F_M$ 表为 $\varinjlim\mathcal{S}_i$，其中 $\mathcal{S}_i$ 是 $F_M$ 的闭子概形，都含零截面 $0\in F_M(A)$。记 $N_i$ 为 $\mathcal{S}_i$ 的余切层%
\SourcePage{324}%
到 $0:\operatorname{Spec}A\hookrightarrow\mathcal{S}_i$ 的限制。则函子 (355) 同构于 (354)，故由 7.12.6(ii) $M$ 是平坦 Mittag-Leffler 模。\hfill$\square$
% [Tutor Note: 原书 p.323 此处 "closed subshemes" 系 "closed subschemes" 之误；底本如此，译文照常作「闭子概形」。]

\emph{注.} 若 $M$ 是任意平坦 $A$-模，则它是有限生成投射 $A$-模 $M_i$ 的有向族的归纳极限，故 $F_M=\varinjlim F_{M_i}$ 是广义意义下的归纳概形（态射 $F_{M_i}\to F_{M_j}$ 不必是闭嵌入）。容易看出：若 $F_M$ 是广义意义下的归纳概形，则 $M$ 平坦。

\bdp{7.12.15} \emph{命题.} 设 $(N_i)_{i\in I}$ 是由有向集 $I$ 参数化的有限生成 $A$-模射影系统，其所有转移映射 $N_j\to N_i$（$j\ge i$）都满。置 $\mathbb{A}(N_i):=\operatorname{Spec}\operatorname{Sym}(N_i)$，$\mathcal{S}:=\varinjlim_i\mathbb{A}(N_i)$。归纳概形 $\mathcal{S}$ 在 $A$ 上形式光滑，当且仅当 $\mathcal{S}$ 同构于 7.12.14 中对应某个平坦 Mittag-Leffler 模 $M$ 的归纳概形 $F_M$。

\emph{证明.} $\mathcal{S}$ 形式光滑，当且仅当函子 (355) 正合（把形式光滑的定义用于形如 $A\oplus J$ 的 $A$-代数，其中 $A\cdot J\subset J$、$J^2=0$）。再用 7.12.6(iii)。\hfill$\square$

\bdp{7.12.16} \emph{命题.} 设 $M$ 是平坦 Mittag-Leffler 模，$F_M$ 是 7.12.14 中的归纳概形。下列条件等价：
\begin{enumerate}[label=(\roman*),leftmargin=2.2em]
\item 对应 $F_M$ 的投射代数（见 7.11.2(i)）是拓扑代数；
\item $M$ 是 [RG] 意义下的严格 Mittag-Leffler 模。
\end{enumerate}

据 [RG] p.74，称 $M$ 为\emph{严格 Mittag-Leffler 模}，若对每个 $f:F\to M$（$F\in\mathcal{C}$）存在 $u:F\to G$（$G\in\mathcal{C}$），使 $f=gu$ 且 $u=hf$，其中 $g:G\to M$、$h:M\to G$（回忆 $\mathcal{C}$ 是有限表现模的范畴）。若 $M=\varinjlim M_i$，$M_i\in\mathcal{C}$，而 $u_{ij}:M_i\to M_j$、$u_i:M_i\to M$ 是典范映射，则 $M$ 严格 Mittag-Leffler 当且仅当：对每个 $i$ 存在 $j\ge i$ 使 $u_{ij}=\varphi_{ij}u_j$，其中 $\varphi_{ij}:M\to M_j$。显然，投射模都是严格 Mittag-Leffler 模，%
\SourcePage{325}%
而严格 Mittag-Leffler 模都是 Mittag-Leffler 模。反之一般不真（见 7.12.24）。
% [Tutor Note: 原书 p.324 此处 "strictly" 误拼作「stritly」；底本如此，译文无法保留拼写，照常译作「严格」。]

\emph{证明.} 把 $M$ 表为 $\varinjlim P_i$，其中 $P_i$ 有限生成投射。置 $N_i:=\operatorname{Im}(P^{*}_j\to P^{*}_i)$，$j$ 充分大。考虑下述条件：
\begin{enumerate}[label=(\alph*),leftmargin=2.2em]
\item 映射 $\varphi_i:\varprojlim_r\operatorname{Sym}(N_r)\to\operatorname{Sym}(N_i)$ 满；
\item $\operatorname{Im}\varphi_i\supset N_i$ 对每个 $i$ 成立；
\item 映射 $\varprojlim_r N_r\to N_i$ 对每个 $i$ 满；
\item 对每个 $i$ 存在 $j\ge i$，使 $\operatorname{Hom}(M,A)$ 与 $\operatorname{Hom}(P_j,A)$ 在 $\operatorname{Hom}(P_i,A)$ 中的像相等。
\end{enumerate}
显然 (i)$\Leftrightarrow$(a)$\Leftrightarrow$(b)$\Leftrightarrow$(c)$\Leftrightarrow$(d)。对 $i\le j$ 考虑映射 $u_{ij}:P_i\to P_j$ 与 $u_i:P_i\to M$。为证 (d)$\Leftrightarrow$(ii)，只须证：$\operatorname{Hom}(M,A)$ 与 $\operatorname{Hom}(P_j,A)$ 在 $\operatorname{Hom}(P_i,A)$ 中的像相等，当且仅当 $u_{ij}=\varphi u_j$，其中 $\varphi:M\to P_j$。证「仅当」：$\operatorname{Hom}(M,P_j)$ 与 $\operatorname{Hom}(P_j,P_j)$ 在 $\operatorname{Hom}(P_i,P_j)$ 中的像相等，故 $\mathrm{id}\in\operatorname{Hom}(P_j,P_j)$ 在 $\operatorname{Hom}(P_i,P_j)$ 中的像，是某个 $\varphi\in\operatorname{Hom}(M,P_j)$ 的像。\hfill$\square$

\bdp{7.12.17} 转入形式光滑仿射 $\aleph_0$-归纳概形的结构之前，先讨论 7.11.1 的形式概形定义与 Grothendieck 的定义（见 EGA I）之间的关系。两者即使在仿射情形也不等价。我们意义下的形式仿射概形，是可表为 $\varinjlim\operatorname{Spec}R_\alpha$ 的归纳概形 $X$，其中 $(R_\alpha)$ 是环的射影系统，映射 $u_{\alpha\beta}:R_\beta\to R_\alpha$（$\beta\ge\alpha$）都满，且 $\operatorname{Ker}u_{\alpha\beta}$ 的元素幂零。Grothendieck 则要求可把 $X$ 表为 $\varinjlim\operatorname{Spec}R_\alpha$，使得映射
\begin{equation}\tag{356}
\varprojlim_\beta R_\beta\to R_\alpha
\end{equation}
% 尾部接 p.326（此句在 p.325 末断开，续文为「是满的」一段）。
